\subsection{Kultur}    \label{subsec_2.1.2}
Was ist Kultur? \textcite{OxfordDic} beschreibt es als zweite Schutzhaut des Replikationsprozesses. Wir Menschen sind als Lebewesen darauf ausgerichtet unsere Spezies am Leben zu erhalten. Jene mit guten Genen überleben besser und pflanzen sich vermehrt fort als solche mit benachteiligten Eigenschaften. In seiner Erklärung, was Kultur ist, beschreibt \textcite{OxfordDic} unsere Haut als den ersten Schutzmechanismus für den Genreplikationsprozess. Da wir jedoch soziale und intelligente Wesen sind haben wir ein weiteres Verteidigungssystem implementiert, das sich außerhalb der Grenzen unseres Körpers befindet; die gemeinsamen kulturellen Errungenschaften. Die Kultur schützt uns demzufolge vor dem Aussterben. Religionen haben sich in den vielen, damals stark von einander getrennte, Kulturen sehr früh als kulturelles System entwickelt, um die Gruppe zu erhalten. Wie lassen sich Religion und Kultur von aneinder trennen? Es ist wichtig festzuhalten, dass sie sich gegenseitig beeinflussen können und dies auch tun. Stark säkulere Länder und Kulturen lassen, durch ihre säkuleren Institutionen und Angeboten für viele Menschen Religion weniger notwendig erscheinen \parencite{ReliInKultur}. So, wie \textcite{Durkeim} Religion definiert, ist sie von der Kultur nicht zu trennen, da ihre Überzeugungen und Praktiken kulturell geprägt sind. Vor vielen Jahrhunderten diktierten die Religionsoberhäupter in Europa die Politik. Von diesem Standpunkt aus ist es nachvollziehbar, dass Kultur und Religion sehr eng mit einander verwoben sind. Die enge Verknüpfung zwischen Kultur und Religion ist jedoch nicht allein auf das Christentum und die westliche Welt zu übertragen. Auch für östliche Religionen, wie dem Buddhismus und dem Islam prägen die religiösen Philosophien die Kultur und umgekehrt \parencite{ReligionCultureandCommunication}. Doch ist dies heutzutage noch aktuell? In Europa beispielsweise schreitet die Sekalarisierung, die Trennung von Kirche und Staat immer weiter voran. In der zweiten Hälfte des 20. Jahrhunderts stieg die Individualisierung der Gesellschaft an und zugleich verloren jüdisch-christliche Institutionen ihre Autorität. Dies hat jedoch der Religion nicht die wichtige Rolle entzogen, in schwierigen und schmerzhaften Zeiten eine Quelle der Kraft und des Haltes zu sein \parencite{ReligionCultureandCommunication}. Obwohl Religion und Kultur interagieren, heißt es nicht, dass sie das gleiche abbilden, oder das gleiche aussagen. Keine Religion der Welt umfasst eine ganze Gesellschaft und keine Gesellschaft richtet sich allein nach religiösen Vorschriften und Praktiken. Es ist sinnvoll, die beiden als miteinander verbunden zu betrachten und ihre Wechselwirkungen zu analysieren \parencite{McGlinchey}.

\subsubsection{Operationalisierung von Kultur}    \label{subsubsec_2.1.2.1}
Wenn Kultur und Religion nicht von einander trennbar ist, wie lässt sich dann Kultur beobachten? \textcite{HofstedeOG} entwickelte sechs Kulturdimensionen, mit denen sich unterschiedliche Kulturen vergleichen und beschreiben lassen. Dafür definierte Hofstede Kultur vereinfacht als \enquote{Software des Geistes} \parencite[S. 209]{HofstedeOG}. Laut ihm setzt sich eine Kultur aus diesen nachfolgenden sechs Dimensionen zusammen und ermöglicht es Kulturen zu vergleichen.

Die \textit{Machtdistanz} beschreibt die Akzeptanz der Machtverteilung. In Kulturen mit hoher Machtdistanz werden starke Hierarchien akzeptiert, in denen beispielsweise Führungskräfte Entscheidungen treffen, die nicht in Frage gestellt werden. Im Gegensatz dazu zeichnen sich Kulturen mit einer niedrigen Machtdistanz durch flache Hierarchien aus, in denen Menschen eher auf gleicher Augenhöhe miteinander interagieren \parencite{HofstedeFocus}.

Die zweite Dimention der \textit{Individualismus} beschreibt die Stärke der zwischenmenschlichen Beziehungen innerhalb einer Kultur. In individualistisch geprägten Gesellschaften sind diese Bindungen eher locker, wobei das Individuum im Mittelpunkt steht. Im Gegensatz dazu steht in kollektivistischen Kulturen das Kollektiv im Vordergrund, wobei Loyalität, Fürsorge und Harmonie wichtige Merkmale sind \parencite{HofstedeFocus}.

\textit{Maskulinität} stellt die dritte Dimension da und beschreibt das Ausmaß, in dem eine Kultur die traditionelle Aufteilung der Geschlechterrollen akzeptiert. In maskulinen Gesellschaften sind diese Rollen klar definiert, während in femininen Kulturen die Rollen flexibler sein können. Darüber hinaus zeichnen sich maskuline Kulturen durch Betonung von Status, Macht und Prestige aus, während feminine Kulturen eher von Fürsorge und Empathie geprägt sind \parencite{HofstedeFocus}.

Die \textit{Unsicherheitsvermeidung} beschreibt das Ausmaß, in dem eine Gesellschaft Risiken scheut und sich gegenüber Veränderungen und Innovationen zurückhaltend verhält. Kulturen mit hoher Unsicherheitsvermeidung bevorzugen bewährte Muster und versuchen, Unsicherheiten durch Regeln und Vorschriften zu minimieren \parencite{HofstedeFocus}.

Die fünfte Dimension \textit{Langzeitorientierung} beschreibt, ob eine Kultur dazu neigt, langfristige Ziele zu planen und zu verfolgen oder eher spontan und kurzfristig handelt und denkt \parencite{HofstedeFocus}.

Die sechste und letzte Dimension \textit{Genuss} zeigt an, inwieweit eine Kultur die Selbstverwirklichung jedes Einzelnen akzeptiert. In Kulturen mit hoher Indulgenz werden auch Randgruppen oder Menschen, die nicht dem Mainstream entsprechen, akzeptiert und die Freiheit wird als wichtiger Wert angesehen. In Kulturen mit niedriger Indulgenz sind die Regeln und Normen hingegen strikter \parencite{HofstedeFocus}.

\textcite{IndiFürStudie} und weitere nutzen die Individualismus-Kollektivismus Dimention um Kulturen in wissenschaftlichen Studien zu vergleichen. Wie bereits in der kurzen Beschreibung der Dimension geschildert, bildet sie die Beziehung des Individuums zur Gesellschaft ab. Steht das Individuum in einer engen Beziehung zu anderen entspricht dies für Kollektivismus. Ist die Beziehung zwischen Individuum und anderen Gesellschaftsmitgliedern loser, dann entspricht dies dem Individualismus \parencite{HofstedeOG}. Individualismus bedeutet aber nicht Egoismus, vielmehr, dass das Individuum eigene Entscheidungen trifft \parencite{Hofstede}. Ein Beispiel dafür ist die USA, in der die Unabhängigkeit und Autonomie des Individuums wertgeschätzt wird \parencite{Utler2020}. Auf der anderen Seite bedeutet Kollektivismus nicht Verbundenheit. Es ist viel mehr, dass jede Personen ihren Platz im Leben kennt, der durch die Gesellschaft bestimmt wird \parencite{Hofstede}. China oder Pakistan sind Beispiele dafür. Dort stellt die Person ihre eigenen Ziele hinter die Ziele der sozialen Gruppe, wie Arbeit oder Familie. Da dies für alle Mitglieder der Gesellschaft gilt, erhalten sie von ihr bedingungslose Unterstützung und Schutz. Auf familiärer Ebene steht die Kernfamilie aus Mutter, Vater, Kind für individualistische und die Großfamilie für kollektivistische Gesellschaften \parencite{Utler2020}. 

% Auswirkung auf die Gesundheit
\subsubsection{Auswirkung auf die Gesundheit}    \label{subsubsec_2.1.2.2}
Die Bedeutung und das Verständnis von Gesundheit, Krankheit und Wohlbefinden sind kulturell geprägt. Dies umfasst, wie Menschen ihre Gesundheit einschätzen, wie sie Symptome interpretieren und wie sie mit gesundheitlichen Herausforderungen umgehen. Diese Verhaltensweisen werden durch kulturelle Normen und Überzeugungen geprägt \parencite{WHO2024, HernandezGibb2020}. Sie beeinflusst die Einstellungen, die Personen zu den verfügbaren Gesundheitsmaßnahmen wie Impfprogrammen, oder Psychotherapie haben und wie sie auf solche Maßnahmen reagieren und kann somit den Erfolg oder das Scheitern dieser Interventionen bestimmen \parencite{WHO2024}.

In vielen Kulturen werden psychische Leiden häufig stigmatisiert oder als körperliche Beschwerden interpretiert, was zu Verzögerungen bei der Suche nach professioneller Hilfe führen kann. Beispielsweise neigen kollektivistische Kulturen, wie China dazu, psychische Probleme zu verschweigen, was zu einer Unterbehandlung führen kann \parencite{HernandezGibb2020}. Da dem Individuum weniger Bedeutung beigemessen wird als beispielsweise der Familie, werden psychische Erkrankungen als Scheitern der Gruppe gesehen. Um das Ansehen der Familie nicht zu beschädigen, werden die Erkrankungen und Beschwerden oft verheimlicht \parencite{KaraszEtAl2019}.

Kulturelle Prägungen spielen im Umgang mit Gesundheitsthemen eine große Rolle. Sie  bestimmen, wie und ob Menschen Zugang zu Gesundheitsdiensten haben. Sie beeinflussen, welche Gesundheitsdienste in Anspruch genommen werden und wie diese genutzt werden \parencite{WHO2024}. So wird im ostasiatischen Raum der Fokus eher auf das Wohl der Gemeinschaft gelegt als auf das Individuum, was sich auf die Gesundheitspraktiken und die Interaktion mit dem Gesundheitssystem auswirkt. In islamischen Gesellschaften kann die Tabuisierung bestimmter Gesundheitsthemen die Gesundheitskompetenz und die Inanspruchnahme von Gesundheitsdiensten negativ beeinflussen \parencite{WielgaKnospe2023}.

Trotz gesundheitlicher Probleme und wirtschaftlicher Härten kann die Kultur einer Gemeinschaft deren Widerstandsfähigkeit stärken und zu einem besseren Wohlbefinden beitragen \parencite{WHO2024}. Des Weiteren können durch kulturelle Praktiken die Bindung zur Gemeinschaft geförtdert werden. Gemeinschaften, die durch gemeinsame kulturelle Traditionen verbunden sind, bieten ihren Mitgliedern ein Gefühl der Zugehörigkeit und emotionaler Unterstützung. Die Gemeinschaft kann in Zeiten von Stress oder Krankheit soziale Unterstützung bieten, was die psychische Gesundheit schützt und stärkt. Kulturelle Praktiken bieten auch physischen Schutz. Dank des Begräbnisrituals der Acholi in Uganda konnten die Ausbreitung der Infektionskrankheit begrenzt werden \parencite{Hewlett2008}. 








% individualistische vs. kollektivistische Kultur
%10 Minutes with Geert Hofstede: https://www.youtube.com/watch?v=zQj1VPNPHlI&t=152s


%\begin{enumerate} [leftmargin=1.25cm] 
%    \item Die \textit{selbstgerichtete Gewalt} beinhaltet suizidales Verhalten sowieGedanken an Suizid und Versuche dessen, sowie die erfolgreiche Vollendung solcher Versuche. Der zweite Bestandteil dieser Gewalt ist die Selbstmisshandlung \parencite{Gewaltarten_WHO}. 
%\end{enumerate}